\documentclass[11pt]{article}

\usepackage{amsmath, amssymb, amsthm}
\usepackage{geometry}
\usepackage{setspace}
\usepackage{lmodern}
\usepackage{microtype}
\usepackage{titlesec}

\geometry{letterpaper, margin=1.2in}
\setstretch{1.15}

% Elegant section formatting
\titleformat{\section}{\large\bfseries}{\thesection.}{0.5em}{}

\begin{document}

\begin{center}
    \LARGE\textbf{The Continuous and the Discrete:} \\[0.2cm]
    \large\textbf{An Analytic Resolution of $G^* = \sqrt{2\pi} \cdot \theta_3(e^{-\pi})^2$}
\end{center}

\vspace{0.4cm}

\noindent There exist rare identities in mathematics that serve as crossroads for entirely disparate fields. The expression $G^* = \sqrt{2\pi} \cdot \theta_3(e^{-\pi})^2$ is one such artifact, sitting at the deep intersection of modular forms, elliptic integrals, continuous probability, and discrete geometry. By unpacking this identity step-by-step, we reveal its exact closed-form value, its algebraic connection to historical constants, and its profound geometric interpretation as a bridge between continuous Euclidean space and the discrete integer lattice.

\section{Evaluating the Jacobi Theta Function}
The core engine of this identity is the Jacobi theta function $\theta_3(q)$, which is defined as the infinite sum of a Gaussian sequence over the integers:
\begin{equation}
    \theta_3(q) = \sum_{n=-\infty}^\infty q^{n^2}
\end{equation}

When we evaluate this at $q = e^{-\pi}$, we are probing the highly symmetric ``self-dual'' point in the complex upper half-plane ($\tau = i$). Because the Gaussian distribution is an eigenfunction of the Fourier transform, evaluating the discrete sum at this precise resonant parameter yields a known exact analytic value governed by the Euler gamma function ($\Gamma$):
\begin{equation}
    \theta_3(e^{-\pi}) = \frac{\pi^{1/4}}{\Gamma(3/4)}
\end{equation}

To render this into a more foundational form, we apply Euler's reflection formula, $\Gamma(z)\Gamma(1-z) = \frac{\pi}{\sin(\pi z)}$. Setting $z = 1/4$:
\begin{align}
    \Gamma(1/4)\Gamma(3/4) &= \frac{\pi}{\sin(\pi/4)} = \pi\sqrt{2} \nonumber \\[0.2cm]
    \implies \Gamma(3/4) &= \frac{\pi\sqrt{2}}{\Gamma(1/4)}
\end{align}

Substituting this reflection back into our evaluated theta function yields the exact amplitude of the one-dimensional discrete Gaussian:
\begin{equation}
    \theta_3(e^{-\pi}) = \frac{\pi^{1/4} \Gamma(1/4)}{\pi\sqrt{2}} = \frac{\Gamma(1/4)}{\sqrt{2}\,\pi^{3/4}}
\end{equation}

\section{Squaring and the Continuous Measure}
Next, the identity squares this discrete theta value. Geometrically, squaring the one-dimensional theta function gives us the partition function of a Gaussian over a two-dimensional square integer lattice ($\mathbb{Z}^2$):
\begin{equation}
    \theta_3(e^{-\pi})^2 = \sum_{m=-\infty}^\infty \sum_{n=-\infty}^\infty e^{-\pi(m^2+n^2)} = \frac{\Gamma(1/4)^2}{2\pi^{3/2}}
\end{equation}

Finally, we multiply this lattice sum by the continuous amplitude $\sqrt{2\pi}$, as prescribed by the identity, to uncover the exact closed algebraic form of $G^*$:
\begin{align}
    G^* &= \sqrt{2\pi} \cdot \frac{\Gamma(1/4)^2}{2\pi^{3/2}} \nonumber \\[0.2cm]
        &= \frac{\sqrt{2}\,\pi^{1/2}\,\Gamma(1/4)^2}{2\pi^{3/2}}
\end{align}

Simplifying the transcendental constants ($\pi^{1/2} / \pi^{3/2} = \pi^{-1}$) yields the true, fundamental nature of the expression:
\begin{equation}
    \boxed{ G^* = \frac{\Gamma(1/4)^2}{\pi\sqrt{2}} \approx 2.9586751\dots }
\end{equation}

\section{A Nexus of Nineteenth-Century Mathematics}
Because $G^*$ is driven fundamentally by the square of the gamma function evaluated at one-quarter, $\Gamma(1/4)^2$, it acts as a mathematical hub that elegantly connects several legendary nineteenth-century constants:

\begin{itemize}
    \item \textbf{Gauss's Constant ($G$):} In 1799, Carl Friedrich Gauss discovered that the reciprocal of the arithmetic-geometric mean ($\operatorname{agm}$) of $1$ and $\sqrt{2}$ yields a fundamental constant of nature:
    \begin{equation*}
        G = \frac{1}{\operatorname{agm}(1, \sqrt{2})} = \frac{\Gamma(1/4)^2}{2\sqrt{2}\,\pi^{3/2}} \approx 0.8346\dots
    \end{equation*}
    Extracting the gamma values from our closed form reveals that the identity is simply a rigorously scaled version of Gauss's constant:
    \begin{equation}
        G^* = 2\sqrt{\pi} \cdot G
    \end{equation}

    \item \textbf{The Lemniscate Constant ($\varpi$):} Just as $\pi$ governs the perimeter of a circle, the constant $\varpi \approx 2.622$ governs the perimeter of the Bernoulli lemniscate (the $\infty$-shaped curve). Because $\varpi = \pi G$, our constant scales directly with the total arc length of the lemniscate:
    \begin{equation}
        G^* = \frac{2}{\sqrt{\pi}} \cdot \varpi
    \end{equation}

    \item \textbf{The Beta Function:} The Euler beta function evaluated at one-quarter is $B(1/4, 1/4) = \frac{\Gamma(1/4)^2}{\sqrt{\pi}}$. This allows us to write the identity elegantly normalized by the continuous standard normal volume:
    \begin{equation}
        G^* = \frac{1}{\sqrt{2\pi}} B(1/4, 1/4)
    \end{equation}

    \item \textbf{Complete Elliptic Integrals:} The complete elliptic integral of the first kind evaluated at the modulus $k=1/\sqrt{2}$ yields $K(1/\sqrt{2}) = \frac{\Gamma(1/4)^2}{4\sqrt{\pi}}$. Thus:
    \begin{equation}
        G^* = \frac{2\sqrt{2}}{\sqrt{\pi}} K(1/\sqrt{2})
    \end{equation}
\end{itemize}

\section{The Continuous $\times$ Discrete Interpretation}
Perhaps the most conceptually poetic way to look at this identity is as a direct, algebraic collision of continuous probability and discrete lattice mathematics.

Isolating the two distinct parts of the original formula reveals their profound geometric origins:
\begin{enumerate}
    \item \textbf{$\sqrt{2\pi}$} is the exact continuous volume (the normalization factor) of the one-dimensional standard Gaussian measure: $\int_{-\infty}^\infty e^{-x^2/2} \, dx$.
    \item \textbf{$\theta_3(e^{-\pi})^2$} is the discrete Gaussian partition function (the heat kernel trace) acting strictly over a rigid, quantized two-dimensional square integer grid, $\mathbb{Z}^2$.
\end{enumerate}

Therefore, the algebraic identity can be conceptually decoupled and rewritten as the tensor product of a continuous space and a discrete boundary:
\begin{equation}
    G^* = \left( \int_{-\infty}^\infty e^{-x^2/2} \, dx \right) \times \left( \sum_{(m,n) \in \mathbb{Z}^2} e^{-\pi(m^2+n^2)} \right)
\end{equation}

\vspace{0.4cm}
\noindent \textbf{Conclusion:} \\
The identity accomplishes something highly remarkable. It takes the continuous probability volume of unbounded Euclidean space ($\sqrt{2\pi}$), marries it to the quantized discrete heat kernel of a two-dimensional integer lattice ($\theta_3(e^{-\pi})^2$), and flawlessly collapses the infinite complexity into a fundamental geometric constant $\left(\frac{\Gamma(1/4)^2}{\pi\sqrt{2}}\right)$. In doing so, it serves as an exact analytic Rosetta Stone---permanently tethering the geometry of the discrete grid to the smooth, unbroken arithmetic-geometric mean discovered by Gauss over two centuries ago.

\end{document}
